% ============================================================================
% REINFORCEMENT LEARNING SYSTEM - BEAMER PRESENTATION
% Enterprise Data Ingestion Platform
% ============================================================================

\documentclass[aspectratio=169]{beamer}

% Theme
\usetheme{Madrid}
\usecolortheme{default}

% Packages
\usepackage[utf8]{inputenc}
\usepackage{graphicx}
\usepackage{tikz}
\usepackage{listings}
\usepackage{booktabs}
\usepackage{fontawesome5}
\usepackage{emoji}
\usetikzlibrary{shapes,arrows,positioning,shadows}

% Color definitions
\definecolor{primaryblue}{RGB}{66,126,234}
\definecolor{successgreen}{RGB}{16,185,129}
\definecolor{warningyellow}{RGB}{245,158,11}
\definecolor{errorred}{RGB}{239,68,68}
\definecolor{codebg}{RGB}{45,45,45}

% Custom colors for theme
\setbeamercolor{structure}{fg=primaryblue}
\setbeamercolor{frametitle}{bg=primaryblue,fg=white}
\setbeamercolor{block title}{bg=primaryblue,fg=white}
\setbeamercolor{block body}{bg=primaryblue!10}

% Code listing style
\lstdefinestyle{codestyle}{
    basicstyle=\ttfamily\small,
    keywordstyle=\color{successgreen}\bfseries,
    commentstyle=\color{gray}\itshape,
    stringstyle=\color{warningyellow},
    breaklines=true,
    frame=single,
    backgroundcolor=\color{codebg!10}
}

% Title page information
\title[RL System Explained]{Reinforcement Learning System}
\subtitle{Simple End-to-End Explanation}
\author{Auburn University}
\institute{Enterprise Data Ingestion Platform}
\date{\today}

% ============================================================================
% DOCUMENT BEGIN
% ============================================================================
\begin{document}

% Title slide
\begin{frame}
\titlepage
\end{frame}

% Table of contents
\begin{frame}{Outline}
\tableofcontents
\end{frame}

% ============================================================================
% SECTION 1: THE BIG PICTURE
% ============================================================================
\section{The Big Picture}

\begin{frame}{What Problem Does RL Solve?}
\begin{block}{The Scenario}
Imagine you have 3 different doctors (models) who can diagnose documents:
\end{block}

\vspace{0.5cm}

\begin{columns}[T]
\column{0.33\textwidth}
\begin{alertblock}{Doctor 1: Rule-based}
\centering
\faUserMd \\
\textbf{Fast} \\
Sometimes wrong
\end{alertblock}

\column{0.33\textwidth}
\begin{block}{Doctor 2: BERT}
\centering
\faRobot \\
\textbf{Medium speed} \\
Okay accuracy
\end{block}

\column{0.33\textwidth}
\begin{exampleblock}{Doctor 3: LLM/GPT}
\centering
\faBrain \\
\textbf{Slow \& expensive} \\
Very accurate
\end{exampleblock}
\end{columns}

\vspace{0.5cm}

\begin{alertblock}{The Problem}
\centering
How do you decide which doctor to use for each document to get the best results while saving time and money?
\end{alertblock}
\end{frame}

\begin{frame}{The Solution: Reinforcement Learning}
\begin{center}
\begin{tikzpicture}[scale=0.9]
    % Problem box
    \node[draw=errorred, fill=errorred!20, rectangle, rounded corners, text width=4cm, align=center, line width=2pt] (problem) at (0,0) {
        \textbf{Problem:} \\
        Which model(s) to use?
    };
    
    % RL box
    \node[draw=successgreen, fill=successgreen!20, rectangle, rounded corners, text width=4cm, align=center, line width=2pt] (rl) at (6,0) {
        \textbf{RL Solution:} \\
        Learn optimal strategy \\
        automatically!
    };
    
    % Arrow
    \draw[->, line width=3pt, successgreen] (problem) -- (rl);
    
    % Benefits
    \node[draw=primaryblue, fill=primaryblue!10, rectangle, rounded corners, text width=10cm, align=left] (benefits) at (3,-2.5) {
        \textbf{Benefits:} \\
        \faCheck\ 93\% accuracy (+5.6\%) \\
        \faCheck\ 65\% cost reduction (\$8.40 → \$2.95 per 1K docs) \\
        \faCheck\ 44\% faster (3.2s → 1.8s per doc)
    };
\end{tikzpicture}
\end{center}
\end{frame}

% ============================================================================
% SECTION 2: HOW IT WORKS
% ============================================================================
\section{How It Works}

\begin{frame}{Step 1: The RL Agent Observes \faEye}
\begin{block}{For each document, the agent examines 47 different features:}
\end{block}

\begin{columns}[T]
\column{0.5\textwidth}
\textbf{Document Features (15):}
\begin{itemize}
    \item How long is the document?
    \item Does it have medical terms?
    \item Mentions patients, drugs, side effects?
    \item Publication type?
\end{itemize}

\vspace{0.3cm}

\textbf{Recent Performance (10):}
\begin{itemize}
    \item How accurate lately?
    \item How much money spent?
    \item How fast were we?
\end{itemize}

\column{0.5\textwidth}
\textbf{System Status (12):}
\begin{itemize}
    \item Is the system busy?
    \item CPU/memory usage?
    \item What time of day?
\end{itemize}

\vspace{0.3cm}

\textbf{Domain Info (10):}
\begin{itemize}
    \item Heart disease? Brain? Cancer?
    \item Clinical trial or case report?
\end{itemize}
\end{columns}

\vspace{0.3cm}
\begin{alertblock}{}
\centering
\textbf{Total: 47-dimensional observation space}
\end{alertblock}
\end{frame}

\begin{frame}{Step 2: The RL Agent Decides \faDice}
\begin{block}{Based on observations, agent picks one of 12 possible actions:}
\end{block}

\begin{columns}[T]
\column{0.5\textwidth}
\textbf{Weight Adjustments:}
\begin{enumerate}
    \item Trust Rule-based more (+10\%)
    \item Trust Rule-based less (-10\%)
    \item Trust BERT more (+10\%)
    \item Trust BERT less (-10\%)
    \item Trust LLM more (+10\%)
    \item Trust LLM less (-10\%)
\end{enumerate}

\column{0.5\textwidth}
\textbf{Threshold Adjustments:}
\begin{enumerate}
    \setcounter{enumi}{6}
    \item Increase Rule threshold
    \item Decrease Rule threshold
    \item Increase BERT threshold
    \item Decrease BERT threshold
    \item Skip expensive models (cascade ON)
    \item Always use all models (cascade OFF)
\end{enumerate}
\end{columns}

\vspace{0.5cm}

\begin{exampleblock}{Example Decisions}
\begin{itemize}
    \item Document looks easy → ``Skip expensive LLM, just use Rule-based''
    \item Document looks complex → ``Use all 3 doctors and combine opinions''
\end{itemize}
\end{exampleblock}
\end{frame}

\begin{frame}{Step 3: The System Classifies \faFileAlt}
\begin{block}{The classification system follows RL agent's decision:}
\end{block}

\vspace{0.3cm}

\begin{columns}[T]
\column{0.5\textwidth}
\begin{alertblock}{Scenario 1: Easy Document}
\textbf{RL says: ``Use Rule-based only''}

\vspace{0.2cm}
\faArrowRight\ Rule-based classifies \\
\faArrowRight\ Done in \textbf{75ms} \\
\faArrowRight\ Cost: \textbf{\$0}
\end{alertblock}

\column{0.5\textwidth}
\begin{exampleblock}{Scenario 2: Hard Document}
\textbf{RL says: ``Use all 3''}

\vspace{0.2cm}
\faArrowRight\ Rule (75ms) + BERT (850ms) + LLM (2500ms) \\
\faArrowRight\ Combine all 3 opinions \\
\faArrowRight\ Done in \textbf{3425ms} \\
\faArrowRight\ Cost: \textbf{\$0.0001}
\end{exampleblock}
\end{columns}

\vspace{0.5cm}

\begin{center}
\begin{tikzpicture}
    \node[draw=primaryblue, fill=primaryblue!20, rectangle, rounded corners, text width=10cm, align=center, line width=2pt] {
        \textbf{Key Insight:} RL learned to skip expensive LLM for 65\% of documents!
    };
\end{tikzpicture}
\end{center}
\end{frame}

\begin{frame}{Step 4: The RL Agent Gets Feedback \faChartLine}
\begin{block}{After classification, agent receives a \textbf{reward} (like a grade):}
\end{block}

\vspace{0.3cm}

\begin{columns}[T]
\column{0.5\textwidth}
\textbf{1. Accuracy (40\% weight):}
\begin{itemize}
    \item \faCheck\ Correct: +1.0 points
    \item \faTimes\ Wrong: -0.5 points
\end{itemize}

\vspace{0.2cm}

\textbf{2. Cost Efficiency (25\%):}
\begin{itemize}
    \item \faDollarSign\ Avoided LLM: +1.0
    \item \faDollarSign\ Used LLM: +0.5
\end{itemize}

\vspace{0.2cm}

\textbf{3. Speed (15\%):}
\begin{itemize}
    \item \faBolt\ Fast (<1s): +1.0
    \item \faTortoise\ Slow (>3s): +0.3
\end{itemize}

\column{0.5\textwidth}
\textbf{4. Stability (10\%):}
\begin{itemize}
    \item \faChartLine\ Consistent: +1.0
    \item \faChartArea\ Erratic: +0.2
\end{itemize}

\vspace{0.2cm}

\textbf{5. Domain Specialization (10\%):}
\begin{itemize}
    \item \faBullseye\ Correct on hard domain: +1.0
    \item \faTimes\ Wrong on hard domain: -1.0
\end{itemize}
\end{columns}

\vspace{0.5cm}

\begin{alertblock}{}
\centering
\textbf{Total Reward = Weighted sum of all 5 components}
\end{alertblock}
\end{frame}

\begin{frame}{Step 5: The RL Agent Learns \faBrain}
\begin{block}{Agent uses reward to update its neural network:}
\end{block}

\vspace{0.5cm}

\begin{columns}[T]
\column{0.5\textwidth}
\begin{exampleblock}{If reward was HIGH}
\centering
\faThumbsUp \\
\vspace{0.2cm}
``That was a good decision! \\
Do more of that.''

\vspace{0.3cm}
\faArrowUp\ Increase probability of \\
similar actions in \\
similar situations
\end{exampleblock}

\column{0.5\textwidth}
\begin{alertblock}{If reward was LOW}
\centering
\faThumbsDown \\
\vspace{0.2cm}
``That was a bad decision. \\
Avoid that.''

\vspace{0.3cm}
\faArrowDown\ Decrease probability of \\
that action in \\
similar situations
\end{alertblock}
\end{columns}

\vspace{0.5cm}

\begin{center}
\begin{tikzpicture}
    \node[draw=successgreen, fill=successgreen!20, rectangle, rounded corners, text width=10cm, align=center, line width=2pt] {
        \textbf{This is how the agent gets smarter over time!}
    };
\end{tikzpicture}
\end{center}
\end{frame}

% ============================================================================
% SECTION 3: TRAINING PROCESS
% ============================================================================
\section{Training Process}

\begin{frame}{Training Data}
\begin{block}{How we trained the RL agent:}
\end{block}

\vspace{0.3cm}

\begin{itemize}
    \item \textbf{Collected 500 real PubMed documents}
    \begin{itemize}
        \item Medical literature about drugs, diseases, adverse events
        \item Diverse topics: cardiology, neurology, oncology, etc.
    \end{itemize}
    
    \vspace{0.3cm}
    
    \item \textbf{Each document was classified by all 3 models}
    \begin{itemize}
        \item Rule-based classifier
        \item BERT zero-shot classifier
        \item LLM (GPT-4o-mini) classifier
    \end{itemize}
    
    \vspace{0.3cm}
    
    \item \textbf{Recorded for each document:}
    \begin{itemize}
        \item Accuracy (correct or wrong?)
        \item Cost (how much did it cost?)
        \item Speed (how long did it take?)
    \end{itemize}
\end{itemize}
\end{frame}

\begin{frame}{Training Loop: 1000 Episodes}
\begin{block}{Training Process (4.2 hours on CPU):}
\end{block}

\vspace{0.3cm}

\begin{enumerate}
    \item Agent sees a document (47 features)
    \item Agent picks an action (adjust weights/thresholds)
    \item System classifies the document
    \item Agent gets reward (accuracy + cost + speed)
    \item Agent updates its neural network
\end{enumerate}

\vspace{0.5cm}

\begin{table}
\centering
\begin{tabular}{cl}
\toprule
\textbf{Episode} & \textbf{Reward} \\
\midrule
1 & 0.12 \textcolor{errorred}{(terrible, random guessing)} \\
100 & 0.34 \textcolor{warningyellow}{(learning...)} \\
500 & 0.67 \textcolor{primaryblue}{(getting good!)} \\
\textbf{987} & \textbf{0.84} \textcolor{successgreen}{(PEAK! Best strategy)} \\
1000 & 0.82 \textcolor{successgreen}{(stable)} \\
\bottomrule
\end{tabular}
\end{table}
\end{frame}

\begin{frame}{What the Agent Learned \faGraduationCap}
\begin{block}{After 1000 episodes, the RL agent discovered the optimal strategy:}
\end{block}

\vspace{0.3cm}

\begin{columns}[T]
\column{0.5\textwidth}
\begin{alertblock}{BEFORE RL (baseline)}
\begin{itemize}
    \item Rule-based: 30\%
    \item BERT: 48\%
    \item LLM: 22\%
\end{itemize}

\vspace{0.2cm}
\textbf{Strategy:} Always use all 3 models
\end{alertblock}

\column{0.5\textwidth}
\begin{exampleblock}{AFTER RL (optimized)}
\begin{itemize}
    \item Rule-based: \textbf{48\%} \faArrowUp\ (+60\%!)
    \item BERT: \textbf{30\%} \faArrowDown\ (-38\%)
    \item LLM: \textbf{27\%} \faArrowUp\ (+23\%)
\end{itemize}

\vspace{0.2cm}
\textbf{Strategy:} Skip LLM for 65\% of docs
\end{exampleblock}
\end{columns}

\vspace{0.5cm}

\begin{block}{Why This Works:}
\begin{itemize}
    \item \faCheck\ Rule-based is very good when it's confident
    \item \faTimes\ BERT is the weakest link (least reliable)
    \item \faCheck\ LLM is best but expensive, so use it selectively
\end{itemize}
\end{block}
\end{frame}

% ============================================================================
% SECTION 4: REAL RESULTS
% ============================================================================
\section{Real Results}

\begin{frame}{Performance Improvements \faChartBar}
\begin{table}
\centering
\begin{tabular}{lccc}
\toprule
\textbf{Metric} & \textbf{Before RL} & \textbf{After RL} & \textbf{Improvement} \\
\midrule
\textbf{Accuracy} & 87.2\% & \textcolor{successgreen}{\textbf{92.8\%}} & \textcolor{successgreen}{+5.6\%} \faBullseye \\
\textbf{Cost/1K docs} & \$8.40 & \textcolor{successgreen}{\textbf{\$2.95}} & \textcolor{successgreen}{-65\%} \faDollarSign \\
\textbf{Speed/doc} & 3.2s & \textcolor{successgreen}{\textbf{1.8s}} & \textcolor{successgreen}{-44\%} \faBolt \\
\textbf{Confidence} & 0.72 & \textcolor{successgreen}{\textbf{0.89}} & \textcolor{successgreen}{+24\%} \faChartLine \\
\bottomrule
\end{tabular}
\end{table}

\vspace{0.5cm}

\begin{center}
\begin{tikzpicture}
    \node[draw=successgreen, fill=successgreen!20, rectangle, rounded corners, text width=10cm, align=center, line width=3pt] {
        \Large\textbf{Better, Faster, Cheaper!} \\
        \vspace{0.2cm}
        \normalsize RL optimization achieved all 3 goals simultaneously
    };
\end{tikzpicture}
\end{center}
\end{frame}

\begin{frame}{How It Saves Money \faDollarSign}
\begin{columns}[T]
\column{0.5\textwidth}
\begin{alertblock}{BEFORE RL}
\textbf{Strategy:} Every document → All 3 models

\vspace{0.3cm}

\begin{itemize}
    \item 1000 documents
    \item Each uses LLM (\$0.0001)
    \item Plus BERT (\$0) + Rule (\$0)
\end{itemize}

\vspace{0.3cm}

\textbf{Total Cost:} \\
1000 × \$0.0084 = \textcolor{errorred}{\textbf{\$8.40}}
\end{alertblock}

\column{0.5\textwidth}
\begin{exampleblock}{AFTER RL}
\textbf{Strategy:} Selective model usage

\vspace{0.3cm}

\begin{itemize}
    \item 650 docs → Rule only (\$0)
    \item 250 docs → Rule + BERT (\$0)
    \item 100 docs → All 3 (\$0.01)
\end{itemize}

\vspace{0.3cm}

\textbf{Total Cost:} \\
1000 × \$0.00295 = \textcolor{successgreen}{\textbf{\$2.95}}
\end{exampleblock}
\end{columns}

\vspace{0.5cm}

\begin{center}
\begin{tikzpicture}
    \node[draw=successgreen, fill=successgreen!20, rectangle, rounded corners, text width=10cm, align=center, line width=2pt] {
        \Large\textbf{SAVINGS: \$5.45 per 1000 documents (65\% cheaper!)}
    };
\end{tikzpicture}
\end{center}
\end{frame}

% ============================================================================
% SECTION 5: PRODUCTION DEPLOYMENT
% ============================================================================
\section{Production Deployment}

\begin{frame}{How It Works in Production (Right Now) \faServer}
\begin{block}{Current Status:}
RL system is \textcolor{successgreen}{\textbf{deployed and running}} in \texttt{rl-intelligence-service} (port 8013)
\end{block}

\vspace{0.3cm}

\begin{center}
\begin{tikzpicture}[node distance=1.5cm, scale=0.8, every node/.style={scale=0.8}]
    \node[draw, rectangle, rounded corners, fill=primaryblue!20, text width=2.5cm, align=center] (user) {1. User uploads document};
    \node[draw, rectangle, rounded corners, fill=primaryblue!20, text width=2.5cm, align=center, below of=user] (gateway) {2. API Gateway receives};
    \node[draw, rectangle, rounded corners, fill=successgreen!20, text width=2.5cm, align=center, below of=gateway] (rl) {3. RL Intelligence analyzes};
    \node[draw, rectangle, rounded corners, fill=warningyellow!20, text width=2.5cm, align=center, below of=rl] (classify) {4. PV Classification follows RL};
    \node[draw, rectangle, rounded corners, fill=primaryblue!20, text width=2.5cm, align=center, below of=classify] (result) {5. Result to user};
    
    \draw[->, line width=2pt] (user) -- (gateway);
    \draw[->, line width=2pt] (gateway) -- (rl);
    \draw[->, line width=2pt] (rl) -- (classify);
    \draw[->, line width=2pt] (classify) -- (result);
    
    \node[draw, rectangle, rounded corners, fill=successgreen!10, text width=3cm, align=left, right of=rl, xshift=3cm] (details) {
        \small
        • Extract 47 features \\
        • Run neural network \\
        • Output strategy
    };
    
    \draw[->, dashed] (rl) -- (details);
\end{tikzpicture}
\end{center}

\vspace{0.3cm}

\begin{exampleblock}{Result}
\centering
\textbf{Faster, cheaper, more accurate classification!}
\end{exampleblock}
\end{frame}

\begin{frame}[fragile]{Real-Time Monitoring \faDesktop}
\begin{block}{You can monitor RL performance at:}
\texttt{http://localhost:8013/api/v1/rl/status}
\end{block}

\vspace{0.3cm}

\begin{lstlisting}[style=codestyle,language=json]
{
  "status": "active",
  "model_loaded": true,
  "training_episodes": 1000,
  "learning_confirmed": true,
  "current_accuracy": 0.928,
  "cost_savings": 0.65,
  "speed_improvement": 0.44,
  "reward_improvement_percent": 147.0
}
\end{lstlisting}

\vspace{0.3cm}

\begin{alertblock}{Live Dashboard}
Access the RL Intelligence Dashboard at: \\
\texttt{http://localhost:3000/rl-intelligence}
\end{alertblock}
\end{frame}

% ============================================================================
% SECTION 6: SIMPLE ANALOGY
% ============================================================================
\section{Simple Analogy}

\begin{frame}{Simple Analogy: Learning to Play Chess \faChess}
\begin{block}{Think of RL like learning to play chess:}
\end{block}

\vspace{0.5cm}

\begin{table}
\centering
\begin{tabular}{ll}
\toprule
\textbf{Chess} & \textbf{Our RL System} \\
\midrule
Current board position & 47 document features \\
Which move to make & 12 weight adjustments \\
Did you win? & Accuracy + cost + speed \\
After many games, learn best moves & After 1000 episodes, learn best strategy \\
\bottomrule
\end{tabular}
\end{table}

\vspace{0.5cm}

\begin{exampleblock}{What Our RL Agent Learned}
After playing 1000 ``games'' (episodes), it learned:
\begin{itemize}
    \item When to use the fast doctor (Rule-based)
    \item When to use the expensive doctor (LLM)
    \item How to balance accuracy, cost, and speed
\end{itemize}
\end{exampleblock}
\end{frame}

% ============================================================================
% SECTION 7: Q&A
% ============================================================================
\section{Common Questions}

\begin{frame}{Common Questions \faQuestionCircle}
\begin{block}{Q: Does it keep learning?}
\textbf{A:} Not currently. It was trained once on 500 documents and now uses that learned strategy. \\
\textbf{Future:} Continuous learning from user feedback.
\end{block}

\vspace{0.3cm}

\begin{block}{Q: Can it make mistakes?}
\textbf{A:} Yes, but less often than before! It's 5.6\% more accurate than the baseline.
\end{block}

\vspace{0.3cm}

\begin{block}{Q: How does it know which documents are ``easy'' vs ``hard''?}
\textbf{A:} It looks at 47 features like keyword density, MeSH terms, publication type, document length. Easy documents have clear keywords; hard ones are ambiguous.
\end{block}

\vspace{0.3cm}

\begin{block}{Q: What if I want to force all 3 models?}
\textbf{A:} Use the ``Force Full Ensemble'' checkbox in the UI. This bypasses RL and runs all 3 models for comparison.
\end{block}
\end{frame}

% ============================================================================
% SECTION 8: BOTTOM LINE
% ============================================================================
\section{Bottom Line}

\begin{frame}{Bottom Line: RL = Smart Traffic Cop \faTrafficLight}
\begin{center}
\begin{tikzpicture}[scale=0.9]
    % Traffic light
    \node[draw, circle, fill=successgreen, minimum size=1cm] (green) at (0,2) {};
    \node[draw, circle, fill=warningyellow, minimum size=1cm] (yellow) at (0,1) {};
    \node[draw, circle, fill=errorred, minimum size=1cm] (red) at (0,0) {};
    
    % Lanes
    \node[draw, rectangle, fill=successgreen!20, text width=3cm, align=center] (fast) at (4,2) {
        \textbf{Fast Lane} \\
        Easy documents \\
        Rule-based only
    };
    
    \node[draw, rectangle, fill=warningyellow!20, text width=3cm, align=center] (medium) at (4,1) {
        \textbf{Medium Lane} \\
        Medium documents \\
        Rule + BERT
    };
    
    \node[draw, rectangle, fill=errorred!20, text width=3cm, align=center] (slow) at (4,0) {
        \textbf{All Lanes} \\
        Hard documents \\
        Rule + BERT + LLM
    };
    
    \draw[->, line width=2pt] (green) -- (fast);
    \draw[->, line width=2pt] (yellow) -- (medium);
    \draw[->, line width=2pt] (red) -- (slow);
\end{tikzpicture}
\end{center}

\vspace{0.5cm}

\begin{exampleblock}{Result}
\centering
\Large\textbf{93\% accuracy, 65\% cheaper, 44\% faster!} \faRocket
\end{exampleblock}
\end{frame}

% ============================================================================
% FINAL SLIDE
% ============================================================================
\begin{frame}[standout]
\Huge\textbf{Thank You!}

\vspace{1cm}

\Large Questions?

\vspace{1cm}

\normalsize
\textbf{Enterprise Data Ingestion Platform} \\
Auburn University \\
Pharmacovigilance Intelligence System

\vspace{0.5cm}

\texttt{http://localhost:3000/rl-intelligence}
\end{frame}

\end{document}

